\section{Elementos sintáticos de uma linguagem de primeira ordem}\label{sec:elementosSintaticosPrimeiraOrdem}
\index{sintaxe}

Nesta seção, trabalharemos um pouco mais com elementos e definições sintáticas de uma linguagem de primeira ordem, como substituições e fechos universais.

\begin{metadefinition}[Variáveis ligadas e livres]\label{mdef:variaveisLigadasELivres}\index{variável!ligada}\index{variável!livre}\index{sentença}
    Dada uma fórmula $\varphi$ e uma ocorrência de um quantificador $Q$ (que pode ser $\forall$ ou $\exists$), dizemos que $Q$ \emph{age} em uma variável $x$ se o próximo símbolo de $\varphi$ após $Q$ for $x$ (ou seja, se o escopo desta ocorrência de $Q$ for da forma $Q x \psi$).

    Dizemos que uma ocorrência de uma variável $x$ em $\varphi$ é \emph{ligada}, ou \emph{quantificada}, se ela estiver no escopo de um quantificador $Q$ que age em $x$.
    Caso contrário, dizemos que a ocorrência de $x$ é \emph{livre}.

    Caso $\varphi$ não possua variáveis livres, dizemos que $\varphi$ é uma \emph{sentença}.
\end{metadefinition}

\begin{example}\label{exa:variaveisLigadasELivres}
    Na fórmula $x_1 = x_2 \wedge \forall \textcolor{red}{x_1} (\textcolor{red}{x_1} \in x_3)$, a segunda e a terceira ocorrência de $x_1$ são ligadas (grafadas em vermelho).
    Todas as outras ocorrências de variáveis são livres.
\end{example}

\begin{example}\label{exa:sentencas}Sobre sentenças:
\begin{itemize}
        \item Se $P$ é um símbolo de predicado de aridade $0$, então $P$ é uma sentença, pois não possui variáveis livres.
    
        \item $\forall x (x = x)$ é uma sentença, pois a única variável que aparece nela é $x$, e todas as ocorrências de $x$ são ligadas.
    
        \item $x = x \wedge \forall x \,(x \in x)$ não é uma sentença, pois as duas primeiras ocorrências de $x$ são livres.
\end{itemize}
\end{example}

A seguinte convenção é muito útil nas mais diversas situações.

\begin{convention}\label{conv:variaveisLivres}
    Seja $\varphi$ uma fórmula.
    A notação $\varphi(x_1, \dots, x_n)$ indica que todas as variáveis com ocorrência livre de $\varphi$ estão entre $x_1, \dots, x_n$.

    Dessa forma, $x_1, \dots, x_n$ podem não ocorrer todas livres em $\varphi$ (não precisam nem mesmo ocorrer em $\varphi$).
    Porém, nenhuma outra variável ocorre livre em $\varphi$.
\end{convention}

Na próxima seção, os axiomas lógicos serão enunciados por meio de esquemas envolvendo fórmulas que podem ter variáveis livres.
Quando for necessário tratar tais fórmulas como sentenças, usaremos fechos universais.
Por isso, o conceito de fecho universal será essencial.

\begin{metadefinition}[Fecho universal]\label{mdef:fechoUniversal}\index{fecho universal}
    Dada uma fórmula $\varphi$, um \emph{fecho universal} de $\varphi$ é uma fórmula do tipo
    \[
        \forall x_1 \dots \forall x_n \varphi,
    \]
    sem variáveis livres, em que $n\geq 0$ e $x_1,\dots,x_n$ são variáveis.

    Em particular, se $\varphi$ já é uma sentença, então $\varphi$ é um fecho universal de si mesma, tomando $n=0$.
    Também permitiremos quantificadores sobre variáveis que não ocorrem livres na fórmula e quantificadores repetidos.
\end{metadefinition}

\begin{example}\label{exa:fechoUniversal}
    Se $\varphi$ é $x_1 = x_2$, então $\forall x_1 \forall x_2 \,(x_1 = x_2)$ é um fecho universal de $\varphi$.
\end{example}

Para trabalharmos com quantificadores, precisaremos da noção de substituição de símbolos em fórmulas, feita como se segue.
Por razões técnicas, também definiremos a substituição de constantes.
Sempre que considerarmos uma constante como alvo de substituição, suas ocorrências serão consideradas livres.

\begin{metadefinition}[Substituição em termos]\label{mdef:substituicaoEmTermos}\index{substituição!em termos}
    Seja $\mathcal L$ um léxico para uma linguagem de primeira ordem, $x$ uma variável ou constante e $t$ um $\mathcal L$-termo.

    Dado um $\mathcal L$-termo $s$, $s[x/t]$ é o termo obtido a partir de $s$ substituindo-se todas as ocorrências de $x$ por $t$.
    Formalmente, define-se $s[x/t]$ por indução sobre a estrutura de $s$, como se segue.
    \begin{itemize}
        \item Se $s$ é $x$, então $s[x/t]$ é $t$.
        \item Se $s$ é uma variável ou constante diferente de $x$, então $s[x/t]$ é $s$.
        \item Se $s$ é $fe_1 \dots e_n$, em que $f$ é um símbolo funcional $n$-ário e $e_1, \dots, e_n$ são termos para os quais $e_1[x/t], \dots, e_n[x/t]$ já foram definidos, então $s[x/t]$ é $fe_1[x/t] \dots e_n[x/t]$.
    \end{itemize}
\end{metadefinition}
Prova-se, por um argumento indutivo simples, que se $s$ é um termo, então $s[x/t]$ é um termo.


Sigamos para alguns exemplos.

\begin{example}\label{exa:substituicaoEmTermos}
    Considere a linguagem da teoria dos anéis do Exemplo~\ref{exa:exemplosLexicos}.
    \begin{itemize}
        \item $x_1[x_1/0]$ é $0$.
        \item $x_1[x_2/0]$ é $x_1$.
        \item $(x_1 + x_2)[x_1/0]$ é $0 + x_2$.
        \item $(x_1 + x_2)[x_2/(x_1 + x_2)]$ é $x_1 + (x_1 + x_2)$.
        \item $(x_1 + 0)[0/(x_1 + x_2)]$ é $x_1 + (x_1 + x_2)$.
    \end{itemize}
\end{example}

A seguir, definiremos a substituição de termos em fórmulas.
Conforme será discutido, é necessária cautela ao se substituir variáveis por termos que contêm variáveis.
Por exemplo, a fórmula $\exists y \, \neg(x = y)$ pode, intuitivamente, ter comportamento diferente conforme o valor atribuído a $x$.
Mas, ao substituir $x$ por $y$, obtemos $\exists y \, \neg(y = y)$, uma fórmula de comportamento lógico completamente diferente.

Optaremos por definir a substituição formal em fórmulas sem nos preocuparmos com este problema, e, a seguir, definir a noção de \emph{substituição válida}.
\begin{metadefinition}[Substituição em fórmulas]\label{mdef:substituicaoEmFormulas}\index{substituição!em fórmulas}
    Seja $\mathcal L$ um léxico para uma linguagem de primeira ordem, $x$ uma variável ou constante e $t$ um $\mathcal L$-termo.

    Dada uma fórmula $\varphi$, $\varphi[x/t]$ é a fórmula obtida por uma substituição literal: substituem-se todas as ocorrências livres de $x$ por $t$ (no caso de constantes, todas as ocorrências), mesmo que variáveis que ocorrem em $t$ passem a ficar no escopo de quantificadores já presentes em $\varphi$.
    
    Formalmente, define-se $\varphi[x/t]$ por indução sobre a estrutura de $\varphi$ pelo algoritmo a seguir.
    \begin{itemize}
        \item Se $\varphi$ é $P e_1 \dots e_n$, em que $P$ é um símbolo de predicado $n$-ário ($n>0$) e $e_1, \dots, e_n$ são termos, então $\varphi[x/t]$ é $P e_1[x/t] \dots e_n[x/t]$.
        \item Se $\varphi$ é $= e_1 e_2$, em que $e_1$ e $e_2$ são termos, então $\varphi[x/t]$ é $= e_1[x/t] e_2[x/t]$.
        \item Se $\varphi$ é um símbolo de predicado 0-ário, então $\varphi[x/t]$ é $\varphi$.
        \item Se $\varphi$ é $\neg \psi$, então $\varphi[x/t]$ é $\neg \psi[x/t]$.
        \item Se $\varphi$ é $\psi \square \theta$, em que $\square$ é $\vee$, $\wedge$, $\rightarrow$ ou $\leftrightarrow$, então $\varphi[x/t]$ é $\psi[x/t] \square \theta[x/t]$.
        \item Se $\varphi$ é $\forall y \psi$ ou $\exists y \psi$, em que $y$ é uma variável diferente de $x$, então $\varphi[x/t]$ é $\forall y \psi[x/t]$ ou $\exists y \psi[x/t]$.
        \item Se $x$ é variável e $\varphi$ é $\forall x \psi$ ou $\exists x \psi$, então $\varphi[x/t]$ é $\varphi$.
    \end{itemize}
\end{metadefinition}

Novamente, por uma indução simples, prova-se que se $\varphi$ é uma fórmula, então $\varphi[x/t]$ é uma fórmula.

Por fim, apresentamos uma definição.

\begin{metadefinition}\label{mdef:substituibilidade}\index{substituição!válida}\index{substituível}
    Seja $\mathcal L$ um léxico para uma linguagem de primeira ordem, $x$ uma variável ou constante e $t$ um $\mathcal L$-termo.
    Dizemos que $x$ é substituível por $t$ em uma fórmula $\varphi$ (ou que a substituição $\varphi[x/t]$ é válida) se nenhuma ocorrência livre de $x$ em $\varphi$ estiver no escopo de um quantificador que age em alguma variável que ocorre em $t$.

    Lembrete: considera-se que toda ocorrência de uma constante é livre.
\end{metadefinition}

Ainda não apresentamos uma noção de verdade ou de dedução, mas, utilizando o conhecimento prévio do leitor, podemos justificar a Metadefinição~\ref{mdef:substituibilidade} com exemplos.

\begin{example}\label{exa:substituicaoInvalida}
    Considere a fórmula $\exists y \, \neg(x = y)$.

    A variável $x$ não é substituível pelo termo $y$ nesta fórmula, pois a ocorrência livre de $x$ está no escopo de um quantificador que age em uma variável que ocorre em $y$ (o próprio $y$).

    Notemos que, ao substituir $x$ por $y$ nessa fórmula, obtemos $\exists y \, \neg(y = y)$.
\end{example}

\begin{example}\label{exa:substituicaoSemVariaveisLivres}
    Considere a fórmula $\varphi$ dada por $\forall x \exists y \, \neg(x = y)$.

    A substituição de $x$ e $y$ por qualquer termo $t$ é válida, pois a fórmula não possui variáveis livres.
    Notemos que, para qualquer termo $t$, $\varphi[x/t]$ é $\forall x \exists y \, \neg(x = y)$ e $\varphi[y/t]$ é $\forall x \exists y \, \neg(x = y)$, ou seja, a fórmula é invariante por tais substituições.
\end{example}

Começaremos com algumas substituições válidas triviais.

\begin{metalemma}
    Seja $\mathcal L$ um léxico para uma linguagem de primeira ordem e $x$ uma variável ou constante.
    Então para todo $\mathcal L$-termo $t$, $t[x/x]=t$, e, para toda fórmula $\varphi$, $x$ é substituível por $x$ em $\varphi$ e $\varphi[x/x]$ é $\varphi$.
\end{metalemma}
\begin{proof}
    Prova-se por indução sobre a estrutura de $t$ que $t[x/x]=t$.

    Se $z$ é uma variável ou constante diferente de $x$, então $z[x/x]$ é $z$.
    Se $z$ for $x$, então $z[x/x]=x=z$.
    Em qualquer caso, na base da indução, temos $t[x/x]=t$.

    Se $t$ é $fe_1 \dots e_n$, em que $f$ é um símbolo funcional $n$-ário e $e_1, \dots, e_n$ são termos para os quais $e_1[x/x]=e_1$, \dots, $e_n[x/x]=e_n$, então $t[x/x]$ é $fe_1[x/x] \dots e_n[x/x]$, que é $fe_1 \dots e_n$.

    Para a segunda afirmação, primeiro note que $x$ é substituível por $x$ em $\varphi$.
    Se $x$ é variável, isso ocorre pois se uma ocorrência de $x$ está no escopo de um quantificador que age em $x$, então tal ocorrência não é livre.
    Se $x$ é constante, isso ocorre pois o termo $x$ não contém variáveis.
    
    Considere uma fórmula atômica $R(e_1, \dots, e_n)$, em que $R$ é um símbolo de predicado $n$-ário e $e_1, \dots, e_n$ são termos.
    Pelo já demonstrado, $R(e_1, \dots, e_n)[x/x]$ é $R(e_1[x/x], \dots, e_n[x/x])$, que é $R(e_1, \dots, e_n)$.
    O caso de fórmulas equacionais é análogo, e o caso de predicados $0$-ários é imediato.

    Se $\varphi$ é $\neg \psi$ e a tese já foi demonstrada para $\psi$, então $\varphi[x/x]$ é $\neg \psi[x/x]$, que é $\neg \psi$.

    Se $\varphi$ é $\psi \square \theta$, em que $\square$ é $\vee$, $\wedge$, $\rightarrow$ ou $\leftrightarrow$, e a tese já foi demonstrada para $\psi$ e para $\theta$, então $\varphi[x/x]$ é $\psi[x/x] \square \theta[x/x]$, que é $\psi \square \theta$.

    Se $\varphi$ é $\forall y \psi$ ou $\exists y \psi$, em que $y$ é uma variável diferente de $x$, e a tese já foi demonstrada para $\psi$, então $\varphi[x/x]$ é $\forall y \psi[x/x]$ ou $\exists y \psi[x/x]$, que é $\forall y \psi$ ou $\exists y \psi$.

    Finalmente, se $\varphi$ é $\forall x \psi$ ou $\exists x \psi$, e a tese já foi demonstrada para $\psi$, então $\varphi[x/x]$ é $\varphi$.
\end{proof}

Também usaremos a seguinte substituição válida trivial.

\begin{metalemma}\label{mlem:substituicaoTrivial}
    Seja $\mathcal L$ um léxico para uma linguagem de primeira ordem e $s$ um $\mathcal L$-termo no qual uma variável ou constante $x$ não ocorre.

    Então para todo $\mathcal L$-termo $t$, $s[x/t]=s$, e, para toda fórmula $\varphi$ em que $x$ não ocorre livre, $x$ é substituível por $t$ em $\varphi$ e $\varphi[x/t]$ é $\varphi$.
\end{metalemma}
\begin{proof}
    Prova-se por indução sobre a estrutura de $s$ que $s[x/t]=s$.
    Se $z$ é uma variável ou constante na qual $x$ não ocorre, então $z\neq x$, de modo que $z[x/t]$ é $z$.

    Se a tese já foi demonstrada para $e_1, \dots, e_n$, e $s$ é $fe_1 \dots e_n$, em que $f$ é um símbolo funcional $n$-ário, então $s[x/t]$ é $fe_1[x/t] \dots e_n[x/t]$, que é $fe_1 \dots e_n$.

    Para a segunda afirmação, primeiro note que $x$ é substituível por $t$ em $\varphi$, pois por hipótese $x$ não possui ocorrências livres em $\varphi$.

    Considere uma fórmula atômica $R(e_1, \dots, e_n)$, em que $R$ é um símbolo de predicado $n$-ário e $e_1, \dots, e_n$ são termos.
    Pelo já demonstrado, $R(e_1, \dots, e_n)[x/t]$ é $R(e_1[x/t], \dots, e_n[x/t])$, que é $R(e_1, \dots, e_n)$.
    O caso de fórmulas equacionais é análogo, e o caso de predicados $0$-ários é imediato.

    Se $\varphi$ é $\neg \psi$ e a tese já foi demonstrada para $\psi$, então $\varphi[x/t]$ é $\neg \psi[x/t]$, que é $\neg \psi$.

    Se $\varphi$ é $\psi \square \theta$, em que $\square$ é $\vee$, $\wedge$, $\rightarrow$ ou $\leftrightarrow$, e a tese já foi demonstrada para $\psi$ e para $\theta$, então $\varphi[x/t]$ é $\psi[x/t] \square \theta[x/t]$, que é $\psi \square \theta$.

    Se $\varphi$ é $\forall y \psi$ ou $\exists y \psi$, em que $y$ é uma variável diferente de $x$, e a tese já foi demonstrada para $\psi$, então $\varphi[x/t]$ é $\forall y \psi[x/t]$ ou $\exists y \psi[x/t]$, que é $\forall y \psi$ ou $\exists y \psi$.

    Finalmente, se $\varphi$ é $\forall x \psi$ ou $\exists x \psi$, e a tese já foi demonstrada para $\psi$, então $\varphi[x/t]$ é $\varphi$.
\end{proof}

O metalema abaixo nos mostra critérios úteis em induções para demonstrar que certas substituições são válidas.

\begin{metalemma}\label{mlem:criteriosSubstituibilidade}
    Seja $\mathcal L$ um léxico para uma linguagem de primeira ordem.
    Sejam $\varphi, \psi$ fórmulas, $x$ uma variável ou constante e $t$ um $\mathcal L$-termo.
    Então:
    \begin{itemize}
        \item Se $\varphi$ é uma fórmula atômica, então $x$ é substituível por $t$ em $\varphi$.
        \item $x$ é substituível por $t$ em $\neg \varphi$ se, e somente se, $x$ é substituível por $t$ em $\varphi$.
        \item $x$ é substituível por $t$ em $\varphi \square \psi$ se, e somente se, $x$ é substituível por $t$ em $\varphi$ e em $\psi$.
        \item Se $x$ é variável, $x$ é substituível por $t$ em $Q x \varphi$, em que $Q$ é $\forall$ ou $\exists$.
        \item $x$ é substituível por $t$ em $Q y \varphi$, em que $Q$ é $\forall$ ou $\exists$ e $y$ é uma variável diferente de $x$, se, e somente se, $x$ é substituível por $t$ em $\varphi$ e ($x$ não ocorre livre em $\varphi$ ou $y$ não ocorre em $t$).
    \end{itemize}
\end{metalemma}
\begin{proof}
    O primeiro item é imediato, uma vez que fórmulas atômicas não possuem quantificadores.

    Para o segundo item, note que os quantificadores de $\neg \varphi$ são exatamente os quantificadores de $\varphi$, com os mesmos escopos.
    
    Para o terceiro item, note que o escopo de um quantificador de $\varphi \square \psi$ é o escopo de um quantificador de $\varphi$ ou o escopo de um quantificador de $\psi$, e vice-versa.

    Para o quarto item, não há ocorrências livres de $x$ em $Qx\varphi$.

    Para o último item, suponha que $x$ é substituível por $t$ em $Qy\varphi$.
    Então nenhuma ocorrência livre de $x$ em $Qy\varphi$ está no escopo de um quantificador que age em alguma variável de $t$.
    Como $x\neq y$, as ocorrências livres de $x$ em $Qy\varphi$ são as mesmas que as ocorrências livres de $x$ em $\varphi$.
    Logo, nenhuma ocorrência livre de $x$ em $\varphi$ está no escopo de um quantificador que age em alguma variável de $t$, ou seja, $x$ é substituível por $t$ em $\varphi$.

    Além disso, se $x$ ocorre livre em $\varphi$, então ocorre livre em $Qy\varphi$, e, portanto, $y$ não ocorre em $t$.

    Reciprocamente, suponha que $x$ é substituível por $t$ em $\varphi$ e que ($x$ não ocorre livre em $\varphi$ ou $y$ não ocorre em $t$).
    Devemos ver que nenhuma ocorrência livre de $x$ em $Qy\varphi$ está no escopo de um quantificador que age em alguma variável de $t$.

    Considere uma ocorrência livre de $x$ em $Qy\varphi$.
    Tal ocorrência é uma ocorrência livre de $x$ em $\varphi$, e, portanto, por hipótese, $y$ não ocorre em $t$.
    Seja $Qw\chi$ o escopo de um quantificador que age em uma variável $w$ que ocorre em $t$, em $Qy\varphi$.
    Como $y\neq w$, tal escopo ocorre em $\varphi$; pela substituibilidade de $x$ por $t$ em $\varphi$, a ocorrência considerada de $x$ não pode estar nesse escopo.
\end{proof}

\begin{metalemma}\label{mlem:substituicaoInversaEmTermos}
    Seja $t$ um $\mathcal L$-termo e sejam $x, y$ variáveis ou constantes.
    Suponha que $y$ não ocorra em $t$.
    Então $x$ não ocorre em $t[x/y]$ e $t[x/y][y/x]$ é $t$.
\end{metalemma}
\begin{proof}
    Prova-se por indução sobre a estrutura de $t$ que $t[x/y][y/x]$ é $t$.
    \begin{itemize}
        \item Se $t$ é uma variável ou constante $z$ diferente de $x$, então $z$ também é diferente de $y$, e portanto $t[x/y][y/x]=z[y/x]=z=t$.
        Note que $x$ não ocorre em $t[x/y]=z$.
        \item Se $t$ é $x$, e $y$ não ocorre em $t$, então $x\neq y$.
        Temos que $t[x/y][y/x]=y[y/x]=x=t$, e $t[x/y]$ é $y$, que não contém $x$.
        \item Se $t$ é $fe_1 \dots e_n$, em que $f$ é um símbolo funcional $n$-ário e $e_1, \dots, e_n$ são termos para os quais $e_1[x/y][y/x]=e_1$, \dots, $e_n[x/y][y/x]=e_n$, então $t[x/y][y/x]$ é $fe_1[x/y][y/x] \dots e_n[x/y][y/x]$, que é $fe_1 \dots e_n$.
        Note que $t[x/y]$ é $fe_1[x/y] \dots e_n[x/y]$, que não contém $x$ porque nenhum dos $e_i[x/y]$ contém $x$.
    \end{itemize}
\end{proof}

\begin{metalemma}\label{mlem:substituicaoInversaEmFormulas}
    Seja $\varphi$ uma $\mathcal L$-fórmula e sejam $x, y$ variáveis ou constantes.
    Suponha que $y$ não ocorra livre em $\varphi$ e que $x$ seja substituível por $y$ em $\varphi$.
    Então $x$ não ocorre livre em $\varphi[x/y]$, $y$ é substituível por $x$ em $\varphi[x/y]$ e $\varphi[x/y][y/x]$ é $\varphi$.
\end{metalemma}
\begin{proof}
    Prova-se por indução sobre a estrutura de $\varphi$.
    Em cada caso, também verificaremos que $y$ é substituível por $x$ em $\varphi[x/y]$.

    \begin{itemize}
        \item Se $\varphi$ é $P e_1 \dots e_n$, em que $P$ é um símbolo de predicado $n$-ário ($n>0$) e $e_1, \dots, e_n$ são termos, então $\varphi[x/y]$ é $P e_1[x/y] \dots e_n[x/y]$.
        Pelo Metalema~\ref{mlem:substituicaoInversaEmTermos}, $x$ não ocorre em cada $e_i[x/y]$.
        Portanto, $x$ não ocorre livre em $\varphi[x/y]$.
        Além disso, como $\varphi[x/y]$ é atômica, $y$ é substituível por $x$ em $\varphi[x/y]$.
        Segue que $\varphi[x/y][y/x]$ é $P e_1[x/y][y/x] \dots e_n[x/y][y/x]$, que é $P e_1 \dots e_n$ pelo Metalema~\ref{mlem:substituicaoInversaEmTermos}.
        \item Se $\varphi$ é $= e_1 e_2$, procedemos como no caso anterior.
        \item Se $\varphi$ é um símbolo de predicado 0-ário $R$, então $\varphi[x/y]$ é $R$, que é $\varphi$.
        \item Se $\varphi$ é $\neg \psi$, então, por hipótese indutiva, $x$ não ocorre livre em $\psi[x/y]$, $y$ é substituível por $x$ em $\psi[x/y]$ e $\psi[x/y][y/x]=\psi$.
        Logo, $x$ não ocorre livre em $\neg\psi[x/y]$, $y$ é substituível por $x$ em $\neg\psi[x/y]$ e $\varphi[x/y][y/x]$ é $\neg \psi[x/y][y/x]$, que é $\neg \psi$.
        \item Se $\varphi$ é $\psi \square \theta$, em que $\square$ é $\vee$, $\wedge$, $\rightarrow$ ou $\leftrightarrow$, então aplicamos a hipótese indutiva a $\psi$ e a $\theta$.
        Logo, $x$ não ocorre livre em $\psi[x/y]\square\theta[x/y]$, $y$ é substituível por $x$ em $\psi[x/y]\square\theta[x/y]$ e $\varphi[x/y][y/x]$ é $\psi[x/y][y/x] \square \theta[x/y][y/x]$, que é $\psi \square \theta$.

        \item Se $\varphi$ é $Q z \psi$, em que $Q$ é $\forall$ ou $\exists$ e $z$ é uma variável diferente de $x$, então $\varphi[x/y]$ é $Q z \psi[x/y]$.
        
        Verifiquemos que nesse caso, com mais detalhes, $y$ é substituível por $x$ em $Q z \psi[x/y]$.
        Se $z=y$, isso é imediato, pois não há ocorrências livres de $y$ em $Qy\psi[x/y]$.
        Se $z\neq y$, então $y$ não ocorre livre em $\psi$.
        Além disso, pelo Metalema~\ref{mlem:criteriosSubstituibilidade}, $x$ é substituível por $y$ em $\psi$.
        Portanto, por hipótese indutiva, $y$ é substituível por $x$ em $\psi[x/y]$.
        Por hipótese, $z$ é diferente de $x$, logo, $z$ não ocorre em $x$.
        Assim, $y$ é substituível por $x$ em $Q z \psi[x/y]$.

        Se $z\neq y$, então, pela hipótese indutiva, $x$ não ocorre livre em $\psi[x/y]$.
        Se $z=y$, então, como veremos abaixo, $x$ não ocorre livre em $\psi$ e, portanto, não ocorre livre em $\psi[x/y]$.
        Logo, em qualquer caso, $x$ não ocorre livre em $Qz\psi[x/y]$.

        Se $z\neq y$, então $\varphi[x/y][y/x]=Qz\psi[x/y][y/x]$, que é $Qz\psi$ pela hipótese indutiva.

        Se $z=y$, então $\varphi[x/y][y/x]$ é $Qz\psi[x/y]$.
        Como $x$ é substituível por $y$ em $Qy\psi$ e $y$ ocorre no termo $y$, segue do Metalema~\ref{mlem:criteriosSubstituibilidade} que $x$ não ocorre livre em $\psi$, e, portanto, $Qz\psi[x/y]$ é $Qz\psi$, ou seja, $\varphi$.

        \item Se $\varphi$ é $Q x \psi$, em que $Q$ é $\forall$ ou $\exists$, então $\varphi[x/y]$ é $Q x \psi$, que é $\varphi$.
        Como $y$ não ocorre livre em $\varphi$, $y$ não ocorre livre em $\varphi[x/y]=\varphi$, e, portanto, $y$ é substituível por $x$ em $\varphi[x/y]=\varphi$.
        Ademais, como $y$ não ocorre livre em $\varphi$, segue que $\varphi[x/y][y/x]=\varphi[y/x]=\varphi$.
    \end{itemize}
\end{proof}

Finalizamos essa seção com um lema técnico que será utilizado mais adiante.

A substituição não é comutativa: considere, por exemplo, a fórmula $\varphi$ dada por $x=y$ na linguagem da teoria dos grupos.
Considere a substituição $[x/x\cdot y]$ e a substituição $[y/x]$.
Temos que $\varphi[x/x\cdot y][y/x]$ é $x\cdot x = x$, enquanto $\varphi[y/x][x/x\cdot y]$ é $x\cdot y = x\cdot y$.

Nesse contexto, vale o seguinte metalema.

\begin{metalemma}\label{mlem:substituicaoComposta}
    Seja $\mathcal L$ um léxico para uma linguagem de primeira ordem.
    Sejam $x$ e $y$ variáveis ou constantes distintas, seja $t$ um $\mathcal L$-termo e seja $u$ um $\mathcal L$-termo no qual $x$ não ocorre.

    Se $r$ é um $\mathcal L$-termo, então $r[x/t][y/u]$ é $r[y/u][x/t[y/u]]$.

    Se $\varphi$ é uma $\mathcal L$-fórmula tal que $x$ é substituível por $t$ em $\varphi$ e $y$ é substituível por $u$ em $\varphi[x/t]$, então $y$ é substituível por $u$ em $\varphi$, $x$ é substituível por $t[y/u]$ em $\varphi[y/u]$ e $\varphi[x/t][y/u]$ é $\varphi[y/u][x/t[y/u]]$.
\end{metalemma}

\begin{proof}
    Primeiro provaremos a afirmação para termos, por indução sobre a estrutura de $r$.

    Se $r$ é uma variável ou constante $z$, temos três casos.

    Se $z=x$, então, como $x\neq y$, $z[x/t][y/u]=t[y/u]=z[y/u][x/t[y/u]]$.

    Se $z=y$, então $z[x/t][y/u]=u$.
    Por outro lado, $z[y/u][x/t[y/u]]=u[x/t[y/u]]=u$, pois, por hipótese, $x$ não ocorre em $u$.
    
    Finalmente, se $z$ é diferente de $x$ e de $y$, então ambos os lados são $z$.

    Se $r$ é $fe_1\dots e_n$, então, pela hipótese indutiva aplicada a $e_1,\dots,e_n$, $e_i[x/t][y/u]=e_i[y/u][x/t[y/u]]$ para cada $i\leq n$.
     Portanto,
    \[
        r[x/t][y/u]
        =
        fe_1[x/t][y/u]\dots e_n[x/t][y/u]
        =
        fe_1[y/u][x/t[y/u]]\dots e_n[y/u][x/t[y/u]],
    \]
    que é exatamente $r[y/u][x/t[y/u]]$.

    Agora provaremos a afirmação para fórmulas, por indução sobre a estrutura de $\varphi$.

    Se $\varphi$ é $P e_1\dots e_n$, em que $P$ é um símbolo de predicado $n$-ário ($n>0$), então a substituibilidade é imediata, pois fórmulas atômicas não possuem quantificadores.
    Além disso, pela parte já provada para termos, $e_i[x/t][y/u]=e_i[y/u][x/t[y/u]]$
    para cada $i\leq n$.
    Logo,
    \[
        \varphi[x/t][y/u]=\varphi[y/u][x/t[y/u]].
    \]

    O caso em que $\varphi$ é $=e_1e_2$ é análogo.
    Se $\varphi$ é um símbolo de predicado 0-ário, nada há a provar.

    Se $\varphi$ é $\neg\psi$, então, pelo Metalema~\ref{mlem:criteriosSubstituibilidade}, as hipóteses implicam que $x$ é substituível por $t$ em $\psi$ e que $y$ é substituível por $u$ em $\psi[x/t]$.
    Pela hipótese indutiva, $y$ é substituível por $u$ em $\psi$, $x$ é substituível por $t[y/u]$ em $\psi[y/u]$ e
    \[
        \psi[x/t][y/u]=\psi[y/u][x/t[y/u]].
    \]
    Aplicando novamente o Metalema~\ref{mlem:criteriosSubstituibilidade}, obtemos as duas substituibilidades desejadas para $\neg\psi$, e a igualdade segue de
    \[
        (\neg\psi)[x/t][y/u]
        =
        \neg\psi[x/t][y/u]
        =
        \neg\psi[y/u][x/t[y/u]]
        =
        (\neg\psi)[y/u][x/t[y/u]].
    \]

    Se $\varphi$ é $\psi\square\theta$, em que $\square$ é $\vee$, $\wedge$, $\rightarrow$ ou $\leftrightarrow$, procedemos da mesma forma, aplicando a hipótese indutiva a $\psi$ e a $\theta$.
    Assim, $y$ é substituível por $u$ em $\psi$ e em $\theta$, e $x$ é substituível por $t[y/u]$ em $\psi[y/u]$ e em $\theta[y/u]$.
    Logo, pelo Metalema~\ref{mlem:criteriosSubstituibilidade}, as substituibilidades desejadas valem para $\psi\square\theta$.
    Além disso,
    \[
        (\psi\square\theta)[x/t][y/u]
        =
        \psi[x/t][y/u]\square\theta[x/t][y/u]
        =
        \psi[y/u][x/t[y/u]]\square\theta[y/u][x/t[y/u]],
    \]
    que é
    \[
        (\psi\square\theta)[y/u][x/t[y/u]].
    \]

    Resta o caso em que $\varphi$ é $Qz\psi$, em que $Q$ é $\forall$ ou $\exists$.

    Se $x$ é variável e $z=x$, então $\varphi[x/t]=\varphi$.
    Como, por hipótese, $y$ é substituível por $u$ em $\varphi[x/t]=\varphi$, a primeira substituibilidade desejada já está provada.
    Além disso, como $x\neq y$, temos
    \[
        \varphi[y/u]=Qx\psi[y/u].
    \]
    Assim, $x$ não ocorre livre em $\varphi[y/u]$, de modo que $x$ é substituível por $t[y/u]$ em $\varphi[y/u]$.
    Por fim, $\varphi[x/t][y/u]=\varphi[y/u]$ e
    \[
        \varphi[y/u][x/t[y/u]]=\varphi[y/u],
    \]
    pois a substituição por $x$ não atravessa o quantificador inicial $Qx$.

    Suponha agora que $z$ é diferente de $x$.
    Então
    \[
        \varphi[x/t]=Qz\psi[x/t].
    \]

    Se $y$ é variável e $z=y$, então $\varphi[y/u]=\varphi$.
    Além disso, $y$ é substituível por $u$ em $\varphi$, pois a substituição por $y$ não atravessa o quantificador inicial $Qy$.
    Da hipótese de que $x$ é substituível por $t$ em $Qy\psi$, segue, pelo Metalema~\ref{mlem:criteriosSubstituibilidade}, que $x$ é substituível por $t$ em $\psi$ e que $x$ não ocorre livre em $\psi$ ou $y$ não ocorre em $t$.

    Se $x$ não ocorre livre em $\psi$, então $x$ também não ocorre livre em $Qy\psi$, e a substituição por $x$ é trivial.
    Se $y$ não ocorre em $t$, então $t[y/u]=t$, e a substituibilidade de $x$ por $t[y/u]$ em $Qy\psi$ é a mesma que a substituibilidade de $x$ por $t$ em $Qy\psi$.
    Em qualquer caso, $x$ é substituível por $t[y/u]$ em $\varphi[y/u]=\varphi$.

    A igualdade também segue nos mesmos dois subcasos.
    Se $x$ não ocorre livre em $\psi$, então ambos os lados são $Qy\psi$.
    Se $y$ não ocorre em $t$, então $t[y/u]=t$, e ambos os lados são $Qy\psi[x/t]$.

    Finalmente, suponha que $z$ é diferente de $x$ e de $y$.
    Como $x$ é substituível por $t$ em $Qz\psi$, o Metalema~\ref{mlem:criteriosSubstituibilidade} implica que $x$ é substituível por $t$ em $\psi$.
    Como $y$ é substituível por $u$ em $Qz\psi[x/t]$, o mesmo metalema implica que $y$ é substituível por $u$ em $\psi[x/t]$.

    Pela hipótese indutiva aplicada a $\psi$, temos que $y$ é substituível por $u$ em $\psi$, que $x$ é substituível por $t[y/u]$ em $\psi[y/u]$ e que
    \[
        \psi[x/t][y/u]=\psi[y/u][x/t[y/u]].
    \]

    Resta verificar as condições referentes ao quantificador inicial $Qz$.
    Da substituibilidade de $y$ por $u$ em $Qz\psi[x/t]$, temos que $y$ não ocorre livre em $\psi[x/t]$ ou $z$ não ocorre em $u$.
    Se $z$ não ocorre em $u$, então a condição para que $y$ seja substituível por $u$ em $Qz\psi$ está satisfeita.
    Se $y$ não ocorre livre em $\psi[x/t]$, então, como $x\neq y$, nenhuma ocorrência livre de $y$ em $\psi$ desaparece ao substituir $x$ por $t$; logo, $y$ não ocorre livre em $\psi$.
    Portanto, em qualquer caso, $y$ é substituível por $u$ em $Qz\psi$.

    Agora vejamos que $x$ é substituível por $t[y/u]$ em $Qz\psi[y/u]$.
    Já sabemos que $x$ é substituível por $t[y/u]$ em $\psi[y/u]$.
    Resta ver que $x$ não ocorre livre em $\psi[y/u]$ ou que $z$ não ocorre em $t[y/u]$.

    Da substituibilidade de $x$ por $t$ em $Qz\psi$, temos que $x$ não ocorre livre em $\psi$ ou $z$ não ocorre em $t$.
    Se $x$ não ocorre livre em $\psi$, então, como $x$ não ocorre em $u$, também $x$ não ocorre livre em $\psi[y/u]$.
    Suponha, então, que $z$ não ocorre em $t$.
    Se $z$ também não ocorre em $u$, então $z$ não ocorre em $t[y/u]$.
    Se $z$ ocorre em $u$, então, pela substituibilidade de $y$ por $u$ em $Qz\psi[x/t]$, temos que $y$ não ocorre livre em $\psi[x/t]$.
    Nesse caso, pela substituibilidade de $x$ por $t$ em $Qz\psi$, uma ocorrência livre de $x$ em $\psi$ produziria uma ocorrência livre de $y$ em $\psi[x/t]$ caso $y$ ocorresse em $t$; concluímos que ou $x$ não ocorre livre em $\psi$, ou $y$ não ocorre em $t$.
    No primeiro caso, $x$ não ocorre livre em $\psi[y/u]$.
    No segundo, como $y$ não ocorre em $t$, temos $t[y/u]=t$, e portanto $z$ não ocorre em $t[y/u]$.

    Assim, $x$ é substituível por $t[y/u]$ em $Qz\psi[y/u]$.

    Por fim, usando a igualdade obtida pela hipótese indutiva,
    \[
        (Qz\psi)[x/t][y/u]
        =
        Qz\psi[x/t][y/u]
        =
        Qz\psi[y/u][x/t[y/u]]
        =
        (Qz\psi)[y/u][x/t[y/u]].
    \]
\end{proof}
\subsection*{Exercícios}

\begin{exercise}\label{exr:variaveisLivresELigadas}
    Determine quais são as ocorrências de variáveis livres e ligadas nas seguintes fórmulas.
    \begin{enumerate}[label=(\alph*)]
        \item $\forall x \, (x = x)$
        \item $\exists x \, (x = y)$
        \item $\forall x \, (x = y) \rightarrow \exists y \, (y = z \wedge x=y)$.
    \end{enumerate}
\end{exercise}

\begin{exercise}\label{exr:substituicaoEmFormulas}
    Considere a linguagem da teoria dos anéis do Exemplo~\ref{exa:exemplosLexicos}.
    Em cada item, determine $\varphi[x/t]$ e se a substituição é válida.
    \begin{enumerate}[label=(\alph*)]
        \item $\varphi$ é $x = x \cdot y$ e $t$ é $y$.
        \item $\varphi$ é $\forall x \, (x = x \cdot y)$ e $t$ é $y$.
        \item $\varphi$ é $\forall x \, (x = x \cdot y \wedge \forall y\, \neg (x = x))$ e $t$ é $y$.
        \item $\varphi$ é $ \neg (x+y =y+x) \wedge \forall y \, (x = x \cdot y)$ e $t$ é $y$.
    \end{enumerate}
\end{exercise}

\begin{exercise}\label{exr:fechosUniversais}
    Determine dois fechos universais distintos para cada uma das fórmulas abaixo.
    \begin{enumerate}[label=(\alph*)]
        \item $x = y$;
        \item $\exists z\,\neg(x = z)$;
        \item $(x = y) \wedge (z = x)$.
    \end{enumerate}
\end{exercise}
